\documentclass[10pt,a4paper]{amsart}
\usepackage[margin=19mm]{geometry}
\usepackage{amsmath,amssymb,mathtools}
\usepackage[scr=rsfs,cal=euler]{mathalfa}
\usepackage{xfrac}
\usepackage[unicode,hidelinks,bookmarks=false]{hyperref}
\newcommand{\ie}{i.e.\ }
\newcommand{\cf}{cf.\ }
\newcommand{\resp}{respectively\ }
\newcommand{\Subsection}{\subsection}
\providecommand{\nobreakdash}{\nobreak\hbox{-}}
\setlength{\emergencystretch}{2em}
\multlinegap0pt
\usepackage{xcolor}
\usepackage{amssymb}
\usepackage{algorithm}\usepackage{algpseudocode}
\usepackage{dsfont}
\newtheorem{theorem}{Theorem}
\newcommand{\E}{\mathbb{E}}
\title{Delightful Exploration}
\author{Ian Osband}
\date{Source: arXiv:2605.13287v1 (CC BY 4.0)}
\begin{document}
\maketitle

\noindent\emph{Source and licence.} Delightful Exploration. Version arXiv:2605.13287v1 (CC BY 4.0), distributed by arXiv under the Creative Commons Attribution 4.0 International licence, http://creativecommons.org/licenses/by/4.0/, as recorded in the arXiv licence metadata for this exact version and shown on the abstract page https://arxiv.org/abs/2605.13287v1. Full licence notice: \texttt{licenses/arxiv-2605.13287-CC-BY-4.0.txt}.\par\smallskip

\noindent\textit{Editorial scope.} Counted unit: the article's theorem thm:tail\_rate (source0.tex lines 772--784). The article's complete original proof of it is delivered with it (source0.tex lines 1366--1434, one \texttt{proof} environment of 69 lines, which the article places away from the statement). No mathematics is written, repaired or re-derived: the statement and the proof are the article's own lines, in the article's own notation, and no retained line is substituted or re-flowed.  No further source-local context is needed: every cross-reference of every retained line resolves inside the two delivered spans.  Nothing was omitted from any retained span, so the package carries no omission record.  No retained line contains a citation, so the wrapper carries no bibliography and no bibliography entry.  No retained line contains a figure environment, an image-inclusion command or drawing code, so no image is delivered and the package declares no asset dependency.  Editorial formatting only, in addition: an AMS \texttt{amsart} preamble that restates the article's own declarations and macros used by the retained lines, namely the environments \texttt{theorem} on separate counters, together with the standard packages those lines need.  The retained environments print in this document as Theorem 1 in order of appearance, whereas the article prints the same items with the numbering of its own full document.  The complete native source of the article as distributed in the arXiv e-print for this exact version is bundled unchanged as \texttt{sources/200491/source0.tex}.
\begin{theorem}[Prior-tail rate for horizon-priced reservation search]
\label{thm:tail_rate}
Assume the upper tail of $F$ is polynomial: for some $\alpha>0$ and constants $0<c_-\le c_+<\infty$,
$c_-y^\alpha \le \Pr(X\ge 1-y) \le c_+y^\alpha$
for all sufficiently small $y$.
Then the reservation policy with $y_T \asymp T^{-1/(\alpha+1)}$ achieves
$R_T = O(T^{\alpha/(\alpha+1)})$.
No policy in the unthrottled infinite-arm discovery model improves this rate in order.
The corresponding DE price obeys $\lambda_T = \Theta(L/T)$.
With constant override rate $\varepsilon$, the same upper-bound argument gives
$R_T = O(\varepsilon^{-1/(\alpha+1)}T^{\alpha/(\alpha+1)})$
and price $\lambda_T = \Theta(L/(\varepsilon T))$.
\end{theorem}

\begin{proof}[Proof of Theorem~\ref{thm:tail_rate}]
For the upper bound, let $\pi_y$ sample fresh arms until observing $X\ge 1-y$, then exploit that arm.
The number of fresh samples before success is geometric with mean $1/p(y)$.
Each exploratory sample incurs regret at most one.
After success, every subsequent round incurs regret at most $y$.
Therefore $R_T(\pi_y)\le 1/p(y)+Ty$.
Using $p(y)\ge c_-y^\alpha$ gives
\[
  R_T(\pi_y)
  \le
  \frac{1}{c_-y^\alpha}+Ty.
\]
Choosing $y_T\asymp T^{-1/(\alpha+1)}$ balances the two terms and yields $R_T=O(T^{\alpha/(\alpha+1)})$.

For the lower bound, fix any policy and any $y>0$.
Let $N_T$ be the number of fresh arms sampled by time $T$, and let $\bar{r}:=1-\E[X]>0$, which is positive under the assumed nontrivial upper-tail model.
Each fresh sample has expected regret $\bar{r}$, so $R_T\ge \bar{r}\,\E[N_T]$.
If $\E[N_T]\ge 1/(4p(y))$, then $R_T\ge \bar{r}/(4p(y))$.
Otherwise $\E[N_T]<1/(4p(y))$.
By a union bound,
\[
  \Pr(\text{some sampled arm has }X\ge 1-y)
  \le p(y)\E[N_T] < 1/4.
\]
With probability at least $3/4$, no sampled arm has value at least $1-y$.
In the revealed-value discovery model, every action played by the policy is either a fresh inspection or a previously inspected arm.
On this event, all such arms have value below $1-y$, so every round incurs regret at least $y$.
Thus $R_T\ge (3/4)Ty$.
Combining the two cases,
\[
  R_T
  \ge
  \min\left\{
    \frac{\bar{r}}{4p(y)},\;
    \frac{3}{4}Ty
  \right\}.
\]
Using $p(y)\le c_+y^\alpha$ and choosing $y\asymp T^{-1/(\alpha+1)}$ gives $R_T=\Omega(T^{\alpha/(\alpha+1)})$.

For the DE price:
\[
  \lambda_T
  =
  L\int_0^{y_T}p(s)\,ds
  =
  \Theta(L\,y_T^{\alpha+1})
  =
  \Theta(L/T).
\]
With constant override rate $\varepsilon$, the expected calendar time to obtain a successful fresh inspection is $1/(\varepsilon p(y))$, so the upper bound becomes
\[
  R_T(y)\le \frac{1}{\varepsilon p(y)}+Ty.
\]
Optimizing gives
\[
  y_T\asymp(\varepsilon T)^{-1/(\alpha+1)},
  \qquad
  R_T=O\!\left(\varepsilon^{-1/(\alpha+1)}T^{\alpha/(\alpha+1)}\right),
\]
and therefore
\[
  \lambda_T
  =
  L\int_0^{y_T}p(s)\,ds
  =
  \Theta(L/(\varepsilon T)).
\]

\end{proof}

\end{document}
